\documentclass{article}

\usepackage[utf8]{inputenc}
\usepackage[T1]{fontenc}
\usepackage{amsmath,amssymb,amsthm}
\usepackage{mathtools}
\usepackage{hyperref}
\usepackage{booktabs}

\makeatletter
\newcommand{\citep}[2][]{%
  \begingroup
  \def\@tmpkey{#2}%
  \ifx\@tmpkey\@empty\else
    (\ifx\@empty#1\@empty\else #1, \fi
     \@nameuse{bibkey@\@tmpkey})%
  \fi
  \endgroup
}
\expandafter\def\csname bibkey@lasser2026micro\endcsname{Microfoundation}
\expandafter\def\csname bibkey@lasser2026paper8b\endcsname{Need~Sufficiency}
\expandafter\def\csname bibkey@lasser2026paper8a\endcsname{Revealed~Sacrifice}
\expandafter\def\csname bibkey@manheim2019\endcsname{Manheim~\&~Garrabrant~2019}
\expandafter\def\csname bibkey@elmhamdi2024\endcsname{El-Mhamdi~\&~Hoang~2024}
\expandafter\def\csname bibkey@majka2025\endcsname{Majka~\&~El-Mhamdi~2025}
\makeatother

\newtheorem{definition}{Definition}
\newtheorem{lemma}{Lemma}
\newtheorem{theorem}{Theorem}
\newtheorem{remark}{Remark}
\newtheorem{assumption}{Assumption}
\newtheorem{proposition}{Proposition}
\newtheorem{corollary}{Corollary}

\newcommand{\HHI}{\mathrm{HHI}}
\newcommand{\Pproxy}{P^{\mathrm{proxy}}}
\newcommand{\Ttruth}{T^{\mathrm{truth}}}
\newcommand{\epsnonres}{\varepsilon_{\mathrm{nonres}}}
\newcommand{\Lip}{\mathrm{Lip}}
\newcommand{\Var}{\mathrm{Var}}
\newcommand{\chisq}{\chi^2}

\title{Concentration-Gap selection theorem, v2: \\
deterministic formulation with $\chi^2$ divergence \\
and norm-level reverse claim}
\author{Working draft for codex review}
\date{}

\begin{document}

\maketitle

\section{Goal and changes from v1}

\subsection{Goal}

Same as v1: discharge a scoped version of the Microfoundation
Concentration-Gap Conjecture along codex's recommended Route D$'$
(algebraic kernel + counterparty interpretation + Route-C transfer).

\subsection{Changes from v1 (codex review responses)}

v1 had 2 P0 + 6 P1 + 1 P2 = 9 findings.  v2 closes them via a
structural reformulation: \textbf{moving from a probabilistic
deviation field to a deterministic one with explicit base population
measure $\mu$}, and \textbf{stating the forward bound via $\chi^2$
divergence} of the trade-flow weighting from the population measure.
This single change resolves the cross-counterparty correlation P0
(no probabilistic decorrelation needed), the random-vs-deterministic
$w$ confusion P1, the type-ambiguous (NCOL) P1, and the
uniform-averaging-on-countable-population P1 all at once.

\begin{enumerate}
\item \textbf{Cross-counterparty correlation P0 fixed.}  v1's
  forward proof was false: (NCOL) controlled weight-deviation
  dependence but not cross-counterparty deviation correlation
  (counterexample: identical centered deviations).  v2's
  deterministic Cauchy-Schwarz argument needs no probabilistic
  decorrelation; the bound holds for any deviation field.

\item \textbf{Lipschitz reverse-claim P0 fixed.}  v1's Theorem 2
  claimed a $g$-level lower bound, but Lipschitz transfer is
  forward-only (constant $g$ counterexample).  v2's reverse claim
  is restricted to the $\|\Pproxy - \Ttruth\|$ level only; the
  $g$-level lower bound is dropped.

\item \textbf{(NCOL) sharpened P1.}  v1's $\|\mathrm{Cov}(w,
  \Delta)\|_{\mathrm{op}} \leq \eta$ was type-ambiguous.  v2 uses
  $\chi^2(w \| \mu)$ which is a single scalar quantity with a
  sharp definition.

\item \textbf{Random-vs-deterministic $w$ P1 fixed.}  v2 treats
  both $w$ and $\Delta$ deterministically; randomness only enters
  if explicitly added (e.g., to model adversary selection
  uncertainty).

\item \textbf{Base population measure P1 added.}  v2 introduces
  a probability measure $\mu$ on the counterparty space
  $\mathcal{C}$; (REP) and (DISP) are stated relative to $\mu$.
  Trade-flow weights $w$ are operational quantities distinct from
  $\mu$.

\item \textbf{Reverse HHI implication P1 fixed.}  v2 uses
  $\max_c w_c \geq \HHI(w)$ (the correct direction; v1 had a
  weaker statement).

\item \textbf{Directional cancellation P1 fixed.}  v2's reverse
  direction uses directional cancellation
  $\langle u_d, \sum_{c \neq d} w_c \Delta_c\rangle \geq -\beta$,
  not the cruder norm bound.

\item \textbf{Embedding/normalization P1 added.}  v2 makes the
  embedding $\phi: \mathcal{V}_\mathrm{utility} \to \mathcal{V}$
  explicit in the counterparty interpretation.

\item \textbf{Coalition closure P1 added.}  v2 makes coalition
  partitioning an explicit pre-theorem audit step; HHI is
  computed over audit-induced coalition equivalence classes.

\item \textbf{Polarization-identity P2 fixed.}  v2 uses
  ``Cauchy-Schwarz'' (the cleaner argument) instead of the
  inappropriate ``parallelogram identity'' phrase.
\end{enumerate}

\section{The algebraic weighted-selection kernel (deterministic)}

\subsection{Setup}

Let $\mathcal{V}$ be a real Hilbert space with norm $\|\cdot\|$
(in applications: capability-poset measures, or its tangent
space).  Let $\mathcal{C}$ be a measurable space (in applications:
the set of counterparties, possibly after coalition partitioning).
Let $\mu$ be a fixed reference probability measure on
$\mathcal{C}$ (the \textbf{base population measure}; in
applications, the auditor-defined ``fair'' or ``natural''
distribution over admissible counterparties).

Let $\Delta: \mathcal{C} \to \mathcal{V}$ be a deterministic
\textbf{deviation field}, $c \mapsto \Delta_c$, $\mu$-integrable.

A \textbf{weighting} $w$ is a probability measure on $\mathcal{C}$
absolutely continuous w.r.t.\ $\mu$, with Radon-Nikodym derivative
$r := \frac{dw}{d\mu}$ satisfying $\mathbb{E}_\mu[r] = 1$.

The \textbf{population mean} of the deviation field is
\begin{equation}
\label{eq:popmean}
  \bar\Delta \;:=\; \int_\mathcal{C} \Delta_c \,d\mu(c) \;=\;
  \mathbb{E}_\mu[\Delta].
\end{equation}

The \textbf{weighted aggregate distortion} is
\begin{equation}
\label{eq:aggregate-distortion-v2}
  \Delta(w) \;:=\; \int_\mathcal{C} \Delta_c \,dw(c)
  \;=\; \mathbb{E}_w[\Delta] \;=\; \mathbb{E}_\mu[r \cdot \Delta].
\end{equation}

\subsection{Two scalar quantities: HHI and $\chi^2$ divergence}

\begin{definition}[$\chi^2$ divergence]
\label{def:chisq}
The \emph{$\chi^2$ divergence} of $w$ from $\mu$ is
\begin{equation}
\label{eq:chisq}
  \chisq(w \,\|\, \mu) \;:=\; \mathbb{E}_\mu[(r - 1)^2]
  \;=\; \mathrm{Var}_\mu(r) \;\geq\; 0.
\end{equation}
$\chisq(w \|\mu) = 0$ iff $w = \mu$ (no selection); larger values
indicate stronger selection of $w$ away from the base measure.
\end{definition}

\begin{definition}[Herfindahl index]
\label{def:hhi-v2}
For discrete $\mathcal{C}$, the \emph{Herfindahl index} of $w$ is
$\HHI(w) := \sum_c w_c^2 = \|w\|_2^2$.

When $\mu$ is the uniform measure on a finite $\mathcal{C}$ of size
$N$, the two divergences are related by
\begin{equation}
\label{eq:chisq-hhi}
  \chisq(w \,\|\, \mu) \;=\; N \cdot \HHI(w) - 1.
\end{equation}
For non-uniform $\mu$, $\chisq(w \|\mu)$ generalizes $\HHI(w)$: it
is the appropriate concentration measure when the population is not
uniformly weighted.
\end{definition}

\subsection{Structural conditions}

\begin{definition}[Representativeness, (REP)]
\label{def:rep-v2}
The deviation field is \emph{$\rho_\mathrm{rep}$-representative}
relative to base measure $\mu$ if
\begin{equation}
  \|\bar\Delta\| \;\leq\; \rho_\mathrm{rep}.
\end{equation}
\end{definition}

\begin{definition}[Bounded dispersion, (DISP)]
\label{def:disp-v2}
The deviation field has \emph{$\sigma$-bounded dispersion}
relative to base measure $\mu$ if
\begin{equation}
  \int_\mathcal{C} \|\Delta_c - \bar\Delta\|^2 \,d\mu(c) \;\leq\;
  \sigma^2.
\end{equation}
\end{definition}

\begin{definition}[Coalition closure, (COAL)]
\label{def:coal-v2}
The counterparty space $\mathcal{C}$ is \emph{coalition-closed}
under audit if all counterparties whose actions or trade flows are
correlated above a deployment-class threshold are partitioned into
a single equivalence class (a coalition), and $\Delta$ is defined
on coalitions rather than individual counterparties.

In particular, $\Delta_c$ for a coalition $c$ is the
coalition-aggregated deviation, and $w_c$ is the
coalition-aggregated weight.
\end{definition}

(COAL) is an audit precondition rather than a probabilistic
condition: the auditor enumerates correlations on the verification
ledger (repeated-beneficiary linkage, governance-vote alignment,
etc.) and partitions counterparties accordingly before the
algebraic theorem is applied.  Latent (undetected) coalition
risk is treated separately; see Remark~\ref{rem:latent-coalitions}.

\subsection{Main forward theorem}

\begin{theorem}[Algebraic weighted-selection, forward direction]
\label{thm:forward-v2}
Let $\Delta$ be a deviation field on $(\mathcal{C}, \mu)$ in a
Hilbert space $\mathcal{V}$, satisfying (REP) and (DISP).  Let $w$
be a weighting on $\mathcal{C}$ with $\chisq(w \|\mu) < \infty$.
Then
\begin{equation}
\label{eq:forward-bound-v2}
  \|\Delta(w)\| \;\leq\; \rho_\mathrm{rep} \;+\;
  \sigma \cdot \sqrt{\chisq(w \,\|\, \mu)}.
\end{equation}
\end{theorem}

\begin{proof}
Decompose $\Delta(w)$ around the population mean:
\[
  \Delta(w) - \bar\Delta \;=\; \mathbb{E}_\mu[r \cdot \Delta]
  - \mathbb{E}_\mu[\Delta] \;=\; \mathbb{E}_\mu[(r - 1) \cdot \Delta]
  \;=\; \mathbb{E}_\mu[(r - 1)(\Delta - \bar\Delta)],
\]
where the last equality uses $\mathbb{E}_\mu[r - 1] = 0$.

We bound the Hilbert-valued mean via duality.  For any unit vector
$u \in \mathcal{V}$ with $\|u\| = 1$:
\[
  \big\langle u, \,\mathbb{E}_\mu[(r - 1)(\Delta - \bar\Delta)] \big\rangle
  \;=\; \mathbb{E}_\mu\big[(r - 1) \cdot \langle u, \Delta - \bar\Delta\rangle\big].
\]
By Cauchy-Schwarz on $L^2(\mu)$ applied to the scalar product of
$(r - 1)$ and $\langle u, \Delta - \bar\Delta\rangle$:
\[
  \big|\mathbb{E}_\mu[(r-1) \langle u, \Delta - \bar\Delta\rangle]\big|
  \;\leq\; \sqrt{\mathbb{E}_\mu[(r-1)^2]} \cdot
  \sqrt{\mathbb{E}_\mu[\langle u, \Delta - \bar\Delta\rangle^2]}.
\]
Since $|\langle u, \Delta - \bar\Delta\rangle| \leq \|\Delta -
\bar\Delta\|$, the second factor is bounded by
$\sqrt{\mathbb{E}_\mu[\|\Delta - \bar\Delta\|^2]} \leq \sigma$.
Taking the supremum over unit $u$ recovers the Hilbert norm:
\[
  \|\mathbb{E}_\mu[(r-1)(\Delta - \bar\Delta)]\|
  \;=\; \sup_{\|u\|=1} \big\langle u, \mathbb{E}_\mu[(r-1)(\Delta - \bar\Delta)]\big\rangle
  \;\leq\; \sqrt{\chisq(w \|\mu)} \cdot \sigma.
\]
Therefore $\|\Delta(w) - \bar\Delta\| \leq \sigma \sqrt{\chisq(w
\|\mu)}$.  Adding the representativeness bound:
\[
  \|\Delta(w)\| \;\leq\; \|\bar\Delta\| + \|\Delta(w) - \bar\Delta\|
  \;\leq\; \rho_\mathrm{rep} + \sigma \sqrt{\chisq(w \|\mu)}.
\]
\end{proof}

\begin{remark}[Why the cross-correlation problem dissolves]
\label{rem:no-cross-corr}
v1's proof attempted to bound
$\mathbb{E}\|\sum_c w_c (\Delta_c - \bar\Delta)\|^2$ via
$\sum_{c,c'} \mathbb{E}[w_c w_{c'}] \mathbb{E}\langle \Delta_c -
\bar\Delta, \Delta_{c'} - \bar\Delta\rangle$, which required a
cross-counterparty decorrelation condition that v1 didn't supply.
The deterministic Cauchy-Schwarz argument in v2 sidesteps this
entirely: the bound depends only on $\chisq(w \|\mu)$ and $\sigma^2$,
both of which are deterministic functionals of the deviation field
and weighting.  Cross-counterparty correlations are absorbed into
the $\sigma^2$ population-variance term automatically.

The codex counterexample (all centered deviations equal a single
random vector $X$) translates to: a deviation field where
$\Delta_c$ is the same for all $c$.  In that case $\sigma^2 = 0$
(no population variance), and the bound gives $\|\Delta(w)\| \leq
\rho_\mathrm{rep}$, with $\Delta(w) = \bar\Delta$ for any $w$.
This is correct: when all counterparties have identical deviations,
selection cannot amplify the distortion.
\end{remark}

\begin{remark}[Latent coalitions]
\label{rem:latent-coalitions}
(COAL) partitions counterparties whose correlations are
\emph{detected} by audit.  Undetected (latent) coalitions among
many small-weight counterparties may produce effective coordination
that the audit's coalition partition does not capture.  The
deployment claim absorbs this risk via an explicit residual term
$\eta_\mathrm{latent}$ that adds to the right-hand side of
\eqref{eq:forward-bound-v2}: deployments where the audit cannot
plausibly bound $\eta_\mathrm{latent}$ are outside the theorem's
scope.  Quantifying $\eta_\mathrm{latent}$ for specific deployments
is operational tooling work, not internal to the algebraic theorem.
\end{remark}

\subsection{Main reverse theorem}

\begin{theorem}[Algebraic weighted-selection, reverse direction]
\label{thm:reverse-v2}
\emph{Scope.}  Throughout this theorem, $\mathcal{C}$ is a finite
or countable atomic measurable space (this is the natural setting
after coalition-closure (COAL) partitions counterparties into
discrete equivalence classes); $w = (w_c)_{c \in \mathcal{C}}$ is
the corresponding atomic weighting.  $\HHI(w)$, $\sum_{c \neq d}$,
and $\arg\max_c w_c$ all refer to this atomic structure.

Let $\Delta$ be a deviation field on $(\mathcal{C}, \mu)$ and $w$
a weighting.  Suppose:
\begin{itemize}
\item (\textbf{Concentration.}) $\HHI(w) \geq H^*$, hence
  $\max_c w_c \geq H^*$ (since $\sum_c w_c^2 \leq \max_c w_c$, the
  Herfindahl is upper-bounded by the maximum weight).  Let $d \in
  \arg\max_c w_c$ be the dominant counterparty, with $w_d \geq
  H^*$.
\item (\textbf{Dominant separation.}) $\|\Delta_d\| \geq \delta$.
\item (\textbf{Directional cancellation.}) Let $u_d :=
  \Delta_d / \|\Delta_d\|$ (unit direction).  The remainder
  cancellation in this direction is bounded:
  \begin{equation}
    \langle u_d, \,\sum_{c \neq d} w_c \Delta_c\rangle \;\geq\;
    -\beta.
  \end{equation}
\end{itemize}

Then
\begin{equation}
\label{eq:reverse-bound-v2}
  \|\Delta(w)\| \;\geq\; H^* \delta - \beta.
\end{equation}
\end{theorem}

\begin{proof}
Project $\Delta(w)$ onto $u_d$:
\begin{align*}
  \langle u_d, \Delta(w)\rangle
  &= w_d \langle u_d, \Delta_d\rangle + \langle u_d,
    \sum_{c \neq d} w_c \Delta_c\rangle \\
  &= w_d \|\Delta_d\| + \langle u_d,
    \sum_{c \neq d} w_c \Delta_c\rangle \\
  &\geq H^* \delta - \beta
\end{align*}
by concentration ($w_d \geq H^*$), separation ($\|\Delta_d\| \geq
\delta$), and directional cancellation.  Then $\|\Delta(w)\| \geq
|\langle u_d, \Delta(w)\rangle| \geq H^* \delta - \beta$, the lower
bound non-trivial when $H^* \delta > \beta$.
\end{proof}

\begin{remark}[Why the reverse needs all three conditions]
\label{rem:reverse-conditions}
Concentration alone does not yield exploitation: a dominant
counterparty whose deviation $\Delta_d$ is small contributes little
distortion.  Separation alone does not yield exploitation: if no
counterparty has dominant share, even large individual deviations
average out.  Directional cancellation is the structural condition
that prevents the remainder counterparties from precisely
counter-balancing the dominant counterparty's contribution.

These three conditions correspond directly to operational audit
hooks: concentration is monitored as paper~10's $I_5$ HHI ceiling;
separation is monitored as utility-divergence between the dominant
counterparty and the population baseline; directional cancellation
is verifiable via remainder-counterparty utility-attestation
patterns.
\end{remark}

\section{Counterparty-selection interpretation}

\subsection{From algebraic kernel to trade-flow setting}

In the GFM trade-flow setting:
\begin{itemize}
\item $\mathcal{C}$: set of counterparties (S1-admissible
  participants under~\citep{lasser2026paper8a}), \emph{after
  coalition closure} per (COAL).
\item $\mu$: base population measure --- the auditor-defined
  ``fair'' or ``natural'' distribution over admissible
  counterparties.  In the simplest case, $\mu$ is uniform; in
  practice, $\mu$ is calibrated to reflect the deployment's
  intended counterparty allocation under representativeness.
\item $\mathcal{V}$: Hilbert space of utility functionals on the
  capability poset (or its tangent space at the welfare-relevant
  truth $W$).
\item Embedding $\phi: \mathcal{V}_\mathrm{utility} \to
  \mathcal{V}$: a normalization map taking each counterparty's
  utility $U_c$ to a comparable Hilbert-space representative.
  The deployment specifies $\phi$ as part of the audit setup.
\item $\Delta_c := \phi(U_c) - \phi(W)$: counterparty $c$'s
  utility deviation from the welfare-relevant truth (in the
  embedded representation).
\item $w_c$: counterparty $c$'s trade-flow weight (paper~10's
  invariant $I_5$ measures $\HHI(w)$ on coalition-closed weights).
\end{itemize}

The aggregate distortion $\Delta(w) = \sum_c w_c \phi(U_c) -
\phi(W)$ is the \textbf{selection-induced proxy-truth deviation}:
how the weighted-trade-flow proxy diverges from the
welfare-relevant truth in the embedded representation.

\subsection{Audit hooks for the structural conditions}

Each of the three structural conditions translates into an
operational audit:

\begin{itemize}
\item \textbf{(REP) representativeness} ($\rho_\mathrm{rep}$):
  \emph{the population's mean utility is bounded close to $W$}.
  Audited by: counterparty-selection-process diversity,
  category-coverage audit (paper~10 Audit~3), attestation-source
  heterogeneity, governance-process diversity.

\item \textbf{(DISP) bounded dispersion} ($\sigma$):
  \emph{individual counterparty utility deviations have bounded
  population variance}.  Audited by: counterparty-level
  utility-disclosure attestations, dispersion estimation from
  observed trade-flow patterns over a calibration window.

\item \textbf{(COAL) coalition closure} ($\eta_\mathrm{latent}$):
  \emph{counterparty correlations above an audit threshold are
  partitioned into coalitions}.  Audited by:
  repeated-beneficiary linkage analysis on the verification
  ledger, governance-vote alignment, substrate-exclusivity
  observability (paper~10 invariant~$I_9$).  Residual undetected
  coordination is bounded by an explicit $\eta_\mathrm{latent}$
  term that the deployment claim absorbs.
\end{itemize}

\section{Route-C transfer: from selection distortion to Goodhart slack}

\subsection{Microfoundation Lipschitz transfer (forward only)}

\citep[Theorem~2]{lasser2026micro}'s Lipschitz transfer states: for
any alignment property $g$ on the operationally active subspace
$P^\mathrm{act}$ with Lipschitz constant $\Lip(g)$,
\begin{equation}
\label{eq:lip-transfer-v2}
  |g(\Ttruth) - g(\Pproxy)| \;\leq\; \Lip(g) \cdot
  \|\Pproxy - \Ttruth\|.
\end{equation}
This is forward-only: $\|\Pproxy - \Ttruth\|$-level bounds carry
to $|g(\Ttruth) - g(\Pproxy)|$-level bounds via $\Lip(g)$, but
\emph{not} the reverse (a constant $g$ has $|g(T) - g(P)| = 0$
regardless of $\|P - T\|$).

\subsection{Identification with selection distortion}

The transfer requires explicit hypotheses on the embedding $\phi$
and on how the selection-induced distortion $\Delta(w)$ relates
to the proxy-truth difference $\Pproxy - \Ttruth$:

\begin{assumption}[Embedding compatibility]
\label{ass:embedding}
The embedding $\phi: \mathcal{V}_\mathrm{utility} \to \mathcal{V}$
is an \emph{affine isometry} on the operationally active subspace
$P^\mathrm{act}$: it preserves the norm, $\|\phi(x) - \phi(y)\| =
\|x - y\|_{\mathcal{V}_\mathrm{utility}}$ for $x, y \in
P^\mathrm{act}$, and is linear up to a fixed translation.

The proxy-truth difference admits the decomposition
\begin{equation}
\label{eq:embedding-decomp}
  \phi(\Pproxy) - \phi(\Ttruth) \;=\; \Delta_\mathrm{audit}(w)
  \;+\; e_\mathrm{latent},
\end{equation}
where $\Delta_\mathrm{audit}(w)$ is the audit-observable selection
distortion (the $\Delta(w)$ of Theorem~\ref{thm:forward-v2}
computed on coalition-closed weights) and $e_\mathrm{latent} \in
\mathcal{V}$ is the latent-coalition residual with
$\|e_\mathrm{latent}\| \leq \eta_\mathrm{latent}$.
\end{assumption}

Under Assumption~\ref{ass:embedding}, isometry gives
$\|\Pproxy - \Ttruth\|_{\mathcal{V}_\mathrm{utility}} = \|\phi(\Pproxy)
- \phi(\Ttruth)\| = \|\Delta_\mathrm{audit}(w) + e_\mathrm{latent}\|$,
and the triangle inequality gives the upper bound used below.  If
$\phi$ is only bi-Lipschitz with constants $\underline{L}, \bar{L}$
(rather than isometric), the same argument applies with bounds
multiplied by $\bar{L}$ and $\underline{L}^{-1}$ respectively; we
state the cleanest form.

\subsection{Theorem 3 (Route-C transfer, forward only)}

\begin{theorem}[Concentration-Gap Selection, Goodhart-slack form]
\label{thm:concgap-v2}
Under (REP), (DISP), (COAL), and Assumption~\ref{ass:embedding}:
\begin{equation}
\label{eq:concgap-bound-v2}
  |g(\Ttruth) - g(\Pproxy)| \;\leq\; \Lip(g) \cdot
  \big(\rho_\mathrm{rep} + \sigma \sqrt{\chisq(w \|\mu)} +
  \eta_\mathrm{latent}\big).
\end{equation}

In particular: $\HHI(w) < H^*$ (paper~10 invariant $I_5$ in force)
combined with uniform $\mu$ on $N$ counterparties bounds
$\chisq(w \|\mu) \leq N H^* - 1$, giving Goodhart slack bounded by
$\Lip(g) \cdot (\rho_\mathrm{rep} + \sigma \sqrt{N H^* - 1} +
\eta_\mathrm{latent})$.

\textbf{Reverse direction (norm-level only).}  Under
Theorem~\ref{thm:reverse-v2}'s concentration + separation +
directional-cancellation conditions:
\begin{equation}
  \|\Pproxy - \Ttruth\| \;\geq\; H^* \delta - \beta.
\end{equation}
A $g$-level lower bound on $|g(\Ttruth) - g(\Pproxy)|$ does
\emph{not} follow from \eqref{eq:lip-transfer-v2} alone; it
requires additional inverse-Lipschitz / non-degeneracy structure
on $g$ that this theorem does not assume.
\end{theorem}

\begin{proof}
Apply \eqref{eq:lip-transfer-v2} and the embedding
identification: $|g(T) - g(P)| \leq \Lip(g) \|P - T\| =
\Lip(g) \|\Delta(w)\| \leq \Lip(g) (\rho_\mathrm{rep} + \sigma
\sqrt{\chisq} + \eta_\mathrm{latent})$ via
Theorem~\ref{thm:forward-v2} extended with the latent-coalition
term.  The reverse follows by Theorem~\ref{thm:reverse-v2} applied
to $\|\Pproxy - \Ttruth\| = \|\Delta(w)\|$.
\end{proof}

\subsection{What this discharges and what it leaves}

\paragraph{Discharged.} A scoped Concentration-Gap theorem:
under (REP) + (DISP) + (COAL) (each operationally auditable), the
proxy-truth Goodhart slack is upper-bounded in terms of $\HHI$ via
$\chi^2$ divergence; under additional separation + directional
cancellation, the proxy-truth norm is lower-bounded.

\paragraph{Replaces SA1.} Paper~10's SA1 (HHI surrogate adequacy:
$\HHI < H^* \Rightarrow$ outside the optimization-pressure regime)
is replaced by the conjunction (REP) + (DISP) + (COAL) plus the
HHI threshold itself.  Each is auditable; SA1's monolithic claim
becomes three concrete checkable conditions.

\paragraph{Does not discharge.} The universal model-free
Concentration-Gap Conjecture: this remains open (and codex's prior
consultation argued it is not provable from current GFM machinery
without a formal optimizer model).  v2 is a scoped theorem with
explicit structural conditions, not the universal claim.

\paragraph{Does not give $g$-level reverse.} The reverse direction
is at the $\|\Pproxy - \Ttruth\|$ level only; no $g$-level lower
bound on Goodhart slack follows from the Lipschitz transfer alone.
For deployment claim purposes, the norm-level bound is sufficient
(it tells operators the regime in which non-trivial proxy-truth
distortion is structurally guaranteed); the $g$-level lower
bound would strengthen the result but requires additional
non-degeneracy structure.

\section{Open questions for codex review}

\begin{enumerate}
\item \textbf{Is the deterministic Cauchy-Schwarz argument
  rigorous?}  Theorem~\ref{thm:forward-v2}'s proof uses the
  $L^2(\mu; \mathcal{V})$ inner product structure.  Is the $L^2$
  setup correct, and does Cauchy-Schwarz apply in the stated form?

\item \textbf{Is the codex counterexample resolved?}
  Remark~\ref{rem:no-cross-corr} argues that the codex
  counterexample translates to a deviation field with
  $\Delta_c$ identical across $c$, hence $\sigma^2 = 0$, hence
  the bound is trivially $\rho_\mathrm{rep}$.  Is this correct?

\item \textbf{Is the latent-coalition treatment adequate?}
  Remark~\ref{rem:latent-coalitions} adds an explicit
  $\eta_\mathrm{latent}$ residual.  Is this the right way to
  handle undetected coordination, or does the theorem need a
  more structural treatment?

\item \textbf{Does the directional-cancellation condition (P1-5
  fix) admit clean operational audit?}  Operationally, an
  auditor would estimate the projection of remainder-counterparty
  utilities onto the dominant's direction.  Is this a feasible
  audit procedure?

\item \textbf{Is the embedding $\phi$ specification adequate?}
  v2 asserts that the deployment specifies $\phi$ as part of the
  audit setup.  Is this enough, or does the theorem need to
  specify $\phi$'s required properties (e.g., metric preservation,
  linearity)?

\item \textbf{Does v2 close all v1 P0/P1 findings adequately?}
  Specifically: P0-1 (cross-counterparty correlation), P0-2
  (Lipschitz reverse), P1-1 (NCOL sharpness), P1-2 ($w$
  random/det), P1-3 (uniform averaging), P1-4 (reverse HHI),
  P1-5 (directional cancellation), P1-6 (embedding), P1-7
  (coalition closure), P2 (terminology) --- each should have a
  specific v2 fix.  Are any inadequately addressed?

\item \textbf{Is v2 ready for integration into a companion paper?}
  If 0 P0 and $\leq$ 2 P1, please declare convergence: ``v2 is
  ready for integration into the companion paper.''  This sets
  the trajectory expectation similar to the C7 cycle.
\end{enumerate}

\section{Status}

v2 attempts to close all v1 P0/P1 findings via the structural
reformulation: deterministic deviation field + base population
measure $\mu$ + $\chi^2$ divergence + restricted reverse claim.

The key structural insight: \textbf{moving from probabilistic to
deterministic resolves the cross-counterparty correlation problem
without needing a separate decorrelation condition} ---
Cauchy-Schwarz in $L^2(\mu; \mathcal{V})$ absorbs cross-correlations
into the $\sigma^2$ population-variance term.

If v2 is sound, the path forward is integration into a companion
paper alongside the C7 closed-form derivation, with paper~10
collapsing its conjecture-dependence to forward references and
dramatically reducing its sensitivity-analysis content.

\end{document}
